\documentclass{article}
\usepackage[T1]{fontenc}
\usepackage{lmodern}
\usepackage{graphicx}
\usepackage{amsmath,amssymb}
\usepackage[a4paper,margin=2cm]{geometry}
\usepackage[absolute,overlay]{textpos}
\setlength{\TPHorizModule}{1mm}
\setlength{\TPVertModule}{1mm}
\usepackage[indonesian]{babel}
\usepackage{hyperref}
\setlength{\emergencystretch}{2em}
% unit: Halaman 14

\begin{document}

% === Halaman 14 (python/native) ===

• Prosedur membuat kunci publik dan kunci privat:

1. Tentukan barisan superincreasing.

2. Kalikan setiap elemen di dalam barisan tersebut dengan n (mod m)


( Modulus m seharusnya angka yang lebih besar daripada jumlah semua 

elemen di dalam barisan, sedangkan pengali n seharusnya tidak 

mempunyai faktor persekutuan dengan m, atau PBB(n, m) = 1 )

3. Hasil perkalian akan menjadi kunci publik sedangkan barisan 

superincreasing semula menjadi kunci privat.

14

\end{document}
